% ============================================================
% Decision Tree Training Pipeline — Code Listing
% CSC491 Textbook, Chapter: Machine Learning
%
% Compile with: xelatex --shell-escape main.tex
%
% Preamble requirements:
%     \usepackage{minted}
%     \usemintedstyle{friendly}   % or tango, vs, xcode
% ============================================================

\begin{listing}[ht]
\begin{minted}[
  linenos,
  frame=lines,
  framesep=3mm,
  fontsize=\small,
  baselinestretch=1.1,
]{python}
from sklearn.tree import DecisionTreeClassifier
from sklearn.preprocessing import StandardScaler


DT_PARAMS = dict(
    max_depth=6,          # limits tree depth to reduce overfitting
    min_samples_leaf=50,  # each leaf must cover at least 50 bars
    class_weight="balanced",  # corrects for imbalanced profit/loss labels
    random_state=42,
)

def train_pipeline(X_train, y_train, X_test, label):
    """Scale features, fit a Decision Tree, and return predictions.

    Returns: (model, scaler, y_prob_test, y_pred_test)
    """
    # Standardise: zero mean, unit variance — required before tree fitting
    scaler   = StandardScaler()
    Xs_train = scaler.fit_transform(X_train)
    Xs_test  = scaler.transform(X_test)       # fit on train only

    dt = DecisionTreeClassifier(**DT_PARAMS)
    dt.fit(Xs_train, y_train)

    y_prob = dt.predict_proba(Xs_test)[:, 1]  # probability of profitable trade
    y_pred = dt.predict(Xs_test)              # binary label: 1 or 0

    print(f"Pipeline {label} trained — depth {dt.get_depth()}")
    return dt, scaler, y_prob, y_pred

# ── Train one model per fractional-differentiation level ─────────────────────
dt_a, sc_a, prob_a, pred_a = train_pipeline(X_train_a, y_train_a, X_test_a, "A (d=0)")
\end{minted}
\caption{%
  Training a decision tree on the financial feature matrix.
  A \texttt{StandardScaler} normalises each feature column before fitting
  (note that \texttt{transform} is called on the test set using parameters
  learned from the training set only, avoiding data leakage).
  Three separate models are trained across the three fractional-differentiation
  levels from the FeatureShootout dataset: $d=0$ (raw prices), $d=1$ (standard
  returns), and $d=x$ (optimal fractional value).
  The \texttt{class\_weight="balanced"} parameter compensates for imbalanced
  profit/loss labels --- a direct application of the pitfall described in
  Section~\ref{sec:pitfalls}.%
}
\label{lst:train-pipeline}
\end{listing}